\documentclass{article}
\usepackage{listings}
\usepackage{amsmath,amsthm}
\usepackage{amssymb}
\usepackage{hyperref}
\usepackage{xurl}


\title{Group projects Programming in Python II}
\author{}
\date{}

\begin{document}
\maketitle

\textbf{The project consists of two parts: (1) the deliverable (code and report) and (2) possibly a short oral exam (online). An oral exam will be scheduled if there is any doubt about the authorship of the submitted code, in particular if there is suspicion that AI tools were used to generate the code. The oral exam will focus on the submitted code.}

\textbf{You may use any resources for the project, but you must not submit code that was generated by AI. Generative AI may be used for finding functions or obtaining ideas, but the submitted code must be written by the team members and you must be able to explain and defend it during the oral exam.}

\textbf{Part of the project grade will be based on code quality (including clarity and comments) and on model performance. You should attempt to optimize model performance and compare different models, selecting the best one.}

\textbf{Deadline: 21/12/2025, 23:59. Submit your *.py or *.ipynb files to the designated folder on the K2 site for the course.}
















\section*{Project 1: Group 2 (Eda	AKKAYA,Ikrame	AMARI), Group 13 (Awa Soraya	BAMBA, Alexis	BOUDIN), 
Group 14 (Awa Soraya	BAMBA,Alexis	BOUDIN), Group 24 (Larissa	DJONKOUE DJANKOU,
Julien	DUPIN), Group 25 (Lea	ASSENHA, Ellischa	BALOURD), Group 30 (Elliot	LOPES, Jinyu	LIN)
}
The task is to predict stock prices based on past stock prices. The data is stored in the file \lstinline$tech_companies_stock_prices.csv$.

The fields are as follows:
\begin{itemize} \item \lstinline$Date$ -- the date of the transaction 
\item \lstinline$Stock_name$ -- the name of the company 
\item \lstinline$Open$, \lstinline$Close$ -- opening and closing prices on the transaction date 
\item \lstinline$Open_i$, \lstinline$Close_i$ -- for $i=1,2,3,4,5$, these represent the opening and closing prices on \lstinline$Date$-$i$ 
\end{itemize}

For each company, predict the share price based on its past share prices over the previous $k$ dates. Specifically: 
\begin{itemize} 
\item Try to predict both the opening and closing prices using various choices of past opening and closing prices. 
\item Explore using past share prices of other companies to predict the share price of a given company. 
\end{itemize}

Use various models (such as linear regression, SVM, decision trees, and neural networks) to predict share prices. Compare the models in terms of prediction accuracy to determine which model performs best.

Finally, select one or two independent variables and use them to build models for predicting share prices. Identify which variables provide the best model for predicting share prices.



\section*{Project 2:  Group 3 (Hafsa	EDDAHNI,
Kévin	KABORE), Group 4 (Yousra	KASMI,
Rali	LAHLOU),  Group 10 (Eva Del Mar	PEREZ,
Valentine	LEJOU), Group 15 (Manassé	LUMBU,Enzo	MATHIAS),
Group 16 (Charly	MENNESSIER,Vitiana	NZAOU), Group 26 (Moetez	BARHOUMI,
Kosseila	BEKHTAOUI), Group 29 (Stella Chanelle	DETIO,Imad	EL ASRI)
}
The task is to predict stock prices based on past stock prices. The data is stored in the file \lstinline$stock_data_for_hw1.csv$.

The fields are as follows: \begin{itemize} \item \lstinline$Date$ -- date of the transaction \item \lstinline$Company$ -- name of the company \item \lstinline$Open$, \lstinline$Close$, \lstinline$High$, \lstinline$Low$ -- opening, closing, highest, and lowest prices on the transaction date \item \lstinline$Open_i$, \lstinline$Close_i$, \lstinline$High_i$, \lstinline$Low_i$ -- for $i=1,2,3,4,5$, these represent the opening, closing, highest, and lowest prices on \lstinline$Date$-$i$ \end{itemize}

For each company, predict the share price based on its past share prices over the previous $k$ dates. Specifically: \begin{itemize} \item Try to predict opening, closing, high, and low prices separately, using different combinations of past opening, closing, high, and low prices. \item Investigate using past share prices of other companies to predict the share price of a given company. \end{itemize}

Use various models (e.g., linear regression, SVM, decision trees, neural networks) for predicting share prices. Compare the models in terms of prediction accuracy to determine which one performs best.

Finally, select one or two independent variables to build models for predicting share prices. Identify which variables yield the best predictive model.





















\section*{Project 3:  Group 5 (Ahmed	OUINEKH,Alexis	VAYSSI
), Group 6 (Bechar	SALEM,Blerina	SHALA), Group 9 (Michel	AFEICHE, Imane	AIT SALAH),
Group 17 (Kaixin	ZHANG,Amal	ZOUAOUI),
Group 18 (Vanessa	SOLARI, Fabrice	TOHON), Group 27 (Tristan	BONAPERA, Mathieu	CAGÉ)
}
 Predict stock prices based on past stock prices, ESG Score, and net income. The CSV file \lstinline$StockESG0.csv$ contains the following columns: \begin{itemize} \item \lstinline$Date$ -- date of transaction \item \lstinline$Company$ -- name of the company \item \lstinline$Open$, \lstinline$Close$ -- opening and closing prices on the transaction date \item \lstinline$Open_i$, \lstinline$Close_i$ -- for $i=1,2,3,4,5$, these represent the opening and closing prices on \lstinline$Date$-$i$ \item \lstinline$Net income$ -- net income of the company \item \lstinline$ESG Score$ -- overall ESG score \item \lstinline$ESG Controversies Score$, \lstinline$Environmental Controversies Count$, \lstinline$Human Rights Controversies$, \lstinline$Business Ethics Controversies$, \lstinline$Wages Working Condition Controversies Count$, \lstinline$Total Renewable Energy To Energy Use in million$, \lstinline$Estimated CO2 Equivalents Emission Total$ -- various ESG-related scores and metrics \end{itemize}

Use past opening or closing prices, net income, and various ESG scores to predict the current opening (or closing) price.

Apply different models, such as linear regression, SVM, decision trees, and neural networks, for predicting the share price. Compare the models in terms of prediction accuracy to determine which performs best.

Finally, select one or two independent variables to build predictive models for share prices. Identify which variables yield the most accurate model. 



 \section*{Project 4:  Group 7 (Axelle	PONCET,Sami	RTEL BENNANI), 
 Group 8 (Yvana	JOUANJAN,Carmelina Grace-Lucia	MBESSO), 
Group 11 (Lucas	GONZALES,Louis	JANIN),
  Group 19 (Alexis	FERRO, Simon	FRADCOURT),
  Group 20 (Hugo ELHAIK, Youssef ELKADY), 
  Group 21 (Asaad	ALSHEIBANI, Ismael	AMINOU), 
Group 28 (Corentin	DAMIAN, Solène	DEBREUIL)
 }
The task is to predict the ESG score based on various other variables, such as: \begin{itemize} \item \lstinline$ESG Controversies Score$ \item \lstinline$Environmental Controversies Count$ \item \lstinline$Human Rights Controversies$ \item \lstinline$Business Ethics Controversies$ \item \lstinline$Wages Working Condition Controversies Count$ \item \lstinline$Total Renewable Energy To Energy Use in million$ \item \lstinline$Estimated CO2 Equivalents Emission Total$ \end{itemize}

The data is stored in the file \lstinline$ESGAllCompanies1.csv$.

Use various models (e.g., linear regression, SVM, decision trees, neural networks) to predict the ESG score. Compare the performance of the models in terms of prediction accuracy. Which model performs the best?

Select one or two independent variables and build models to predict the ESG score. Identify which variables provide the best predictive model for the ESG score. 

\section*{Project 5:  Group 12 (Maiwenn SIMON, Aude NEVO),  Group 22 (Paul GRAVIS ,Pierre YVENOU ), Group 1 (Rim ZAMZAMI,Titouan VAN MEYEL), Group 23 (Gabriel SOYER, Mattis LEFEBURE))}
   The task is to reverse engineer a portfolio management algorithm.
   The historical stock data is stored in the file \lstinline$tech_companies_stock_prices.csv$. The fields are as follows:
   \begin{itemize}
   \item \lstinline$Date$ -- the date of the transaction
   \item \lstinline$Stock_name$ -- the name of the company
   \item \lstinline$Open$, \lstinline$Close$ -- opening and closing prices on the transaction date
   \item \lstinline$Open_i$, \lstinline$Close_i$ -- for $i=1,2,3,4,5$, these represent the opening and closing prices on \lstinline$Date$-$i$
   \end{itemize}

    A portfolio is a collection of stocks and cash. A portfolio management algorithm decides, at the beginning of each trading day, how many stocks to buy or sell from each company,
   based on the current portfolio (numbers of stocks and cash) and the stock prices (opening and closing prices) of the previous days.
   The goal of the portfolio management algorithm is to maximize the value of the portfolio in the long run.


   The historical performance of the portfolio management algorithm can be found in \lstinline$test_portfolio.csv$.
   The rows of this file are indexed by dates of transactions, the columns correspond to various stocks. In addition, there is a
   column \lstinline$cash$ and a column \lstinline$value$. The column \lstinline$value$ represents the value of the portfolio at the end of the given day.
   The values in columns labeled by stocks correspond to the number of stocks bought (if positive) or sold (if negative) at the
   beginning of the day, using the opening prices. For computing the value, the cash in the portfolio and the values of the stocks in the
   portfolio using the closing prices are added.
   It is assumed that at the beginning of the day, the proceeds from selling stocks at the opening prices and the cash reserve
   can be used to buy stock at the opening prices.

   For example, the line
   \begin{lstlisting}
      ,    TSLA, AAPL, AMZN, MSFT, NVDA, META, GOOGL, cash,               value
     2023-10-10,0.0,  4.0,  -2.0, -2.0, 0.0,  0.0,  1.0,   68.65997314452716, 303305.98236083984
   \end{lstlisting}
   means that on \lstinline$2023-10-10$, at the beginning of the day, $4$ more shares of \lstinline$AAPL$ were bought, $2$ shares of \lstinline$AMZN$ and \lstinline$MSFT$ were
   sold, and $1$ more share of \lstinline$GOOGL$ was bought, and \lstinline$68.65997314452716$ was added to the cash reserves. At the end of the day,
   the value of the portfolio was \lstinline$303305.98236083984$, where the closing prices were used for computing the value.
   Note that $4$ is not the number of \lstinline$AAPL$ shares in the portfolio, but the number of \lstinline$AAPL$ shares which were bought in addition to
   the \lstinline$AAPL$ shares already in the portfolio. Likewise, \lstinline$-2.0$ is the number of \lstinline$AMZN$ shares sold from the number of \lstinline$AMZN$ shares in
   the portfolio.

   First, use the historical data from \lstinline$tech_companies_stock_prices.csv$ and \lstinline$test_portfolio.csv$
   to reconstruct the portfolio (numbers of stocks and cash) at the end of each trading date. Second, use various models
   (e.g., linear regression, SVM, decision trees, neural networks) to predict:
   \begin{itemize}
   \item the value of the portfolio at the end of each day (except the first trading day) as a function of the portfolio at the end of the previous day and the opening prices of the current day;
   \item the trading actions (buying/selling stocks) as a function of the trading actions during the previous $k$ days ($k=1,2,3$).
   \end{itemize}

   Keep at least 20\% of the data for testing.

   %Try to propose a better trading strategy than the one implemented in \lstinline$test_portfolio.csv$, and compare its performance
   %on the test data, which was not used for reverse engineering of the trading strategy.

\end{document}
The task is to predict the logarithm of house prices (variable \emph{lprice})d using the following variables: \begin{itemize} \item logarithm of the lot size (variable \emph{llotsize}) \item logarithm of the surface area of the house (variable \emph{lsqrft}) \item number of bedrooms (variable \emph{bdrms}) \item surface area of the house (variable \emph{sqrft}) \item lot size (variable \emph{lotsize}) \end{itemize}

The data is stored in the file \lstinline$hprice1_merged.csv$.

Use various models (such as linear regression, SVM, decision trees, and neural networks) to predict the logarithm of the house price. Compare the performance of the models in terms of prediction accuracy. Which model performs the best?

Select one or two independent variables and build models to predict the logarithm of the price. Identify which variables provide the best model for predicting the logarithm of the price.

Repeat the exercise by replacing the logarithm of the price with the price itself (variable \emph{price}).


\end{document}
